\documentclass[11pt,a4paper]{article}

\usepackage[T1]{fontenc}
\usepackage[utf8]{inputenc}
\usepackage[margin=0.95in]{geometry}
\usepackage{mathpazo}
\usepackage{microtype}
\usepackage{enumitem}
\usepackage{parskip}
\usepackage{titling}
\usepackage[hidelinks]{hyperref}

\setlength{\parindent}{0pt}
\setlength{\parskip}{5pt}

\setlist[itemize]{leftmargin=1.4em, topsep=2pt, itemsep=2pt}
\setlist[enumerate]{leftmargin=1.6em, topsep=2pt, itemsep=2pt}

\setlength{\droptitle}{-3.5em}
\posttitle{\par\end{center}\vspace{-1em}}

\title{A One-Page Diagnostic for AI Errors in Legal Work\\[0.2em]
\large Five questions to ask when an AI-assisted draft looks
plausible but is not quite right.}

\author{\small Lodewijk Bonebakker\thanks{Drafting assistance:
Claude Opus 4.7 (Anthropic). The author is not a practicing
attorney; this is framework, not legal advice.}}

\date{\small April 2026}

\begin{document}
\maketitle
\vspace{-1.5em}

\noindent
After \emph{Mata v.~Avianca} [1] and the Stanford RegLab finding
of 17--33\% hallucination rates even in retrieval-grounded legal
AI tools [2], firms have learned to fear hallucination. But ``the
AI hallucinated'' is a catch-all that covers at least five
distinguishable failures, each with a different cause and a
different remedy. Catching only fabricated cases leaves four other
failure modes unchecked. Walk the ladder bottom-up:

\textbf{1. Structure --- is the output well-formed?}
Bluebooked citations, consistent defined terms, correct filing
format. Easy to catch with linters, citation checkers, and
contract validators. Necessary, not sufficient: a correctly
formatted document can still be wrong about everything else.

\textbf{2. Meaning --- is it internally coherent?}
A contract whose recitals contradict its covenants; a redlined
clause that inverts the risk allocation; a memo whose conclusion
does not follow from the analysis. Detectable by careful reading
without recourse to outside sources. Domain expertise unmasks the
subtle ones.

\textbf{3. Context --- is it appropriate for the situation?}
A demand letter in aggressive tone when the client wants to
preserve the commercial relationship; a research memo surveying
federal law for a state-court matter; a brief written for the
wrong audience. The reviewer can usually name what is wrong; the
fix is to give the model more of the situation, in writing.

\textbf{4. Groundedness --- do the legal claims connect to actual
law?} Fabricated cases, misquoted holdings, overruled authority
cited as good law. The post-\emph{Mata} danger zone, and the one
the profession has learned hardest. Verify any cited authority
that matters; treat AI claims as hypotheses; mandate KeyCite or
Shepard's verification. Verification of groundedness is now table
stakes. When groundedness failures appear consistently in a
specific practice area --- specialty courts, narrow regulatory
domains, emerging statutory areas --- the tool may be
miscalibrated to that niche; practice-area evaluation before
deployment is the structural remedy, not per-matter
verification.

\textbf{5. Environment --- did the prompt encode what you know?}
The model does not know the partner's strategic posture, the local
court's standing orders, the client's risk tolerance, or the
matter's history. The output can be fluent, internally coherent,
properly grounded, and still wrong because the question itself was
underspecified. The lawyer's job is to write the situation into
the prompt; no automated tool can do this for the lawyer.

\textbf{For multi-step analysis.} Walk each material intermediate
step --- standing, jurisdiction, standard of review, element
satisfaction --- not only the final conclusion. A chain where
every individual step passes the ladder is not therefore sound;
a subtle error in an early step propagates through the downstream
reasoning without any single link appearing defective.

\textbf{Bottom line.}
Structure and groundedness are largely tooling problems. Meaning,
context, and environment are where the lawyer's judgment is most
clearly indispensable, and where the continuing duty of competence
and supervision under ABA Formal Opinion 512 [3] most clearly
applies. A firm that wants to reduce groundedness errors should
invest in retrieval and citation auditing. A firm that wants to
reduce environment errors should invest in training, templates,
and reviewer practice. Knowing which rung you are on tells you
which budget line to charge.

\vspace{0.5em}
\hrule
\vspace{0.4em}

{\small
\noindent
[1] \emph{Mata v.~Avianca, Inc.}, No.~22-cv-1461, 2023 WL 4114965
(S.D.N.Y. June 22, 2023).
\quad
[2] V.~Magesh, F.~Surani, M.~Dahl, M.~Suzgun, C.~D.~Manning,
D.~E.~Ho, ``Hallucination-Free? Assessing the Reliability of
Leading AI Legal Research Tools,'' Stanford RegLab / HAI, 2024.
\quad
[3] ABA Formal Opinion 512, \emph{Generative Artificial
Intelligence Tools}, July 29, 2024.
}

\end{document}
